\documentclass[12pt]{article}
\usepackage[utf8]{inputenc}
\usepackage[brazil]{babel}
\usepackage{graphicx}
\usepackage{amsmath}
\usepackage{hyperref}
\usepackage{caption}
\usepackage{geometry}
\usepackage{enumitem}
\usepackage{booktabs}
\usepackage{multirow}
\usepackage{colortbl}
\usepackage{xcolor}
\usepackage{subcaption}
\usepackage{tikz,tikz-3dplot}
\usepackage[linesnumbered,ruled,vlined]{algorithm2e}

\geometry{a4paper, margin=2.5cm}
\setlength{\parskip}{0.3\baselineskip}

\hypersetup{
    colorlinks=true,
    linkcolor=black,
    citecolor=black,
    urlcolor=black
}

%=======================================================
%  INDIQUE AQUI O NOME DO PROBLEMA E A EQUIPE
%=======================================================
\title{Universidade Federal de Juiz de Fora \\
        DCC059 - Teoria dos Grafos \\
        \vspace{6.0cm}
        {\Huge\textbf{Árvore Geradora Mínima\\
                      Capacitada (CMSTP)}}\\
        \vspace{6.0cm}
}
\author{
    Thomás Causin \\
    Caio Cunha \\
    Igor Reis \\
}
\date{Julho de 2026}

%=======================================================
%=======================================================
\begin{document}

\maketitle
\newpage

\tableofcontents
\newpage

%=======================================================
%                 I N T R O D U Ç Ã O
%=======================================================
\section{Introdução}

O presente relatório descreve o projeto de implementação desenvolvido para o \textit{Capacitated Minimum Spanning Tree Problem} (CMSTP), o Problema da Árvore Geradora Mínima Capacitada. O trabalho insere-se no contexto da disciplina de Teoria dos Grafos, onde foram estudados algoritmos exatos para problemas clássicos — como Prim e Kruskal para árvores geradoras mínimas. O CMSTP é uma generalização do problema de árvore geradora mínima que, ao acrescentar uma restrição de capacidade sobre as sub-árvores, torna-se NP-difícil \cite{garey1979}, inviabilizando o uso direto desses algoritmos exatos em instâncias de porte médio e grande.

Foram desenvolvidos três algoritmos construtivos baseados na heurística de Esau--Williams \cite{esauwilliams1966}: um algoritmo guloso determinístico, um algoritmo guloso randomizado (GRASP \cite{feo1995}) e um algoritmo guloso randomizado reativo \cite{prais2000}. Esses algoritmos foram avaliados sobre instâncias clássicas da OR-Library \cite{beasley1990} e os resultados foram comparados com as melhores soluções conhecidas da literatura \cite{ruiz2015}.

O restante do trabalho está assim estruturado: na Seção~\ref{secProblema} o problema é descrito formalmente; a Seção~\ref{secAlgoritmos} detalha as abordagens propostas; a Seção~\ref{secResultados} apresenta os experimentos computacionais e a análise dos resultados; por fim, a Seção~\ref{secConclusoes} traz as conclusões e propostas de trabalhos futuros.


%=======================================================
%     D E S C R I Ç Ã O     D O    P R O B L E M A
%=======================================================
\section{Descrição do Problema}
\label{secProblema}

\subsection{Motivação e modelagem em grafos}

O CMSTP modela um cenário recorrente em redes de telecomunicações: dado um \textbf{concentrador central} (a raiz) e um conjunto de \textbf{terminais} (os demais vértices), deseja-se interligá-los em uma árvore de custo mínimo, respeitando a restrição de que cada ``cabo'' que sai diretamente do concentrador não pode transportar mais do que $Q$ unidades de tráfego.

Formalmente, dado um grafo não-direcionado completo $G = (V, E)$, onde:
\begin{itemize}[noitemsep]
    \item $r \in V$ é o vértice raiz (concentrador), com $|V| = n$;
    \item $c_{ij} \geq 0$ é o custo da aresta entre os vértices $i$ e $j$;
    \item $d_i > 0$ é a demanda do terminal $i$ (com $d_r = 0$);
    \item $Q > 0$ é a capacidade máxima de cada sub-árvore,
\end{itemize}
o objetivo é encontrar uma \textbf{árvore geradora enraizada em $r$} de custo mínimo tal que, ao remover a raiz, cada sub-árvore resultante (denominada \textit{s-tree}) tenha soma de demandas no máximo $Q$:

\begin{eqnarray}
    \min \sum_{e \in E_T} c_e && \label{func_obj} \\
    \text{sujeito a} \quad \sum_{i \in S} d_i \leq Q, && \forall\; \text{s-tree}\; S. \label{restricao_cap}
\end{eqnarray}

A função objetivo~(\ref{func_obj}) minimiza o custo total das arestas da árvore. A restrição~(\ref{restricao_cap}) impõe que a soma das demandas dos terminais de cada s-tree não exceda a capacidade $Q$.

A Figura~\ref{fig:instancia_solucao} ilustra uma instância com 4 terminais e uma solução com duas s-trees.

\begin{figure}[ht!]
  \label{fig:instancia_solucao}
  \centering

  \begin{subfigure}{0.45\textwidth}
    \centering
    \resizebox{0.95\textwidth}{!}{
    % Instância CMSTP de exemplo: raiz 0 (quadrado) + terminais 1-4 (círculos)
% Arestas sólidas: ligações à raiz; tracejadas: arestas internas
\begin{tikzpicture}
\tikzset{
    >=stealth,
    aresta/.style={-, thick},
    arestaCinza/.style={-, thin, dashed, gray},
    vertice/.style={circle, semithick, draw=black, inner sep=0pt, minimum size=0.7cm},
    raiz/.style={rectangle, very thick, draw=black, fill=gray!20,
                 minimum width=0.8cm, minimum height=0.7cm}
}

% Vértices
\node[raiz]    (v0) at (2.5,2)   {\footnotesize \textbf{0}};
\node[vertice] (v1) at (0,  3.8) {\footnotesize 1};
\node[vertice] (v2) at (0,  0.2) {\footnotesize 2};
\node[vertice] (v3) at (5,  3.8) {\footnotesize 3};
\node[vertice] (v4) at (5,  0.2) {\footnotesize 4};

% Ligações à raiz (custo da aresta)
\draw[aresta] (v0) -- node[above left,  font=\footnotesize] {10} (v1);
\draw[aresta] (v0) -- node[below left,  font=\footnotesize] {12} (v2);
\draw[aresta] (v0) -- node[above right, font=\footnotesize] {11} (v3);
\draw[aresta] (v0) -- node[below right, font=\footnotesize] {9}  (v4);

% Arestas internas
\draw[arestaCinza] (v1) -- node[left,   font=\footnotesize] {3} (v2);
\draw[arestaCinza] (v3) -- node[right,  font=\footnotesize] {4} (v4);
\draw[arestaCinza] (v2) -- node[below,  font=\footnotesize] {5} (v4);
\draw[arestaCinza] (v1) -- node[above,  font=\footnotesize] {8} (v3);

% Legenda
\node[font=\footnotesize, anchor=north west] at (6.3, 4.4) {Legenda:};
\draw[aresta]      (6.3, 3.8) -- (7.3, 3.8);
\node[font=\footnotesize, anchor=west] at (7.4, 3.8) {ligação à raiz};
\draw[arestaCinza] (6.3, 3.2) -- (7.3, 3.2);
\node[font=\footnotesize, anchor=west] at (7.4, 3.2) {aresta interna};
\draw[thick] (6.1,2.4) rectangle (9.6,4.6);
\end{tikzpicture}

    }
    \caption{Instância CMSTP: raiz \textbf{0} (retângulo), terminais 1--4 (círculos), $Q = 3$, demandas unitárias.}
    \label{fig:instancia}
  \end{subfigure}
  \hfill
  \begin{subfigure}{0.45\textwidth}
    \centering
    \resizebox{0.95\textwidth}{!}{
    % Solução após Esau-Williams: s-tree {1,2} e s-tree {3,4}
\begin{tikzpicture}
\tikzset{
    >=stealth,
    gate/.style={-, thick, blue!70!black},
    interna/.style={-, thick, teal!80!black},
    vertice/.style={circle, semithick, draw=black, inner sep=0pt, minimum size=0.7cm},
    raiz/.style={rectangle, very thick, draw=black, fill=gray!20,
                 minimum width=0.8cm, minimum height=0.7cm}
}
\node[raiz]    (v0) at (2.5, 2)   {\footnotesize \textbf{0}};
\node[vertice] (v1) at (0,   3.5) {\footnotesize 1};
\node[vertice] (v2) at (-1,  1.5) {\footnotesize 2};
\node[vertice] (v3) at (5,   3.5) {\footnotesize 3};
\node[vertice] (v4) at (6,   1.5) {\footnotesize 4};
% Gates (ligações de cada s-tree à raiz)
\draw[gate] (v0) -- node[above, font=\footnotesize] {10} (v1);
\draw[gate] (v0) -- node[above, font=\footnotesize] {9}  (v4);
% Arestas internas (economias realizadas)
\draw[interna] (v1) -- node[left,  font=\footnotesize] {3} (v2);
\draw[interna] (v4) -- node[right, font=\footnotesize] {4} (v3);
% Rótulos das s-trees
\node[font=\scriptsize, color=blue!70!black] at (-0.8, 4.1) {s-tree $\{1,2\}$};
\node[font=\scriptsize, color=blue!70!black] at (5.8,  4.1) {s-tree $\{3,4\}$};
\node[font=\small, align=center] at (2.5,-0.7) {Custo $= 10+3+9+4 = 26$};
\end{tikzpicture}

    }
    \caption{Solução Esau--Williams: s-tree $\{1,2\}$ e s-tree $\{3,4\}$, cada uma com demanda $2 \leq Q$.}
    \label{fig:solucao}
  \end{subfigure}

  \caption{Instância e solução do exemplo CMSTP ($n=5$, $Q=3$).}
\end{figure}

\subsection{Complexidade e necessidade de heurísticas}

A AGM sem restrições é resolvida em tempo polinomial pelos algoritmos de Kruskal \cite{kruskal1956} e Prim. Porém, ao acrescentar a restrição de capacidade~(\ref{restricao_cap}), o CMSTP torna-se NP-difícil por redução ao Problema da Mochila \cite{garey1979}. Isso significa que, salvo $P=NP$, não existe algoritmo eficiente capaz de garantir a solução ótima para instâncias grandes, o que justifica o uso de heurísticas.

\subsection{Por que o Kruskal não resolve o CMSTP?}

Uma pergunta natural é: não bastaria aplicar o algoritmo de Kruskal \cite{kruskal1956} e depois verificar a capacidade? A resposta é não, por dois motivos fundamentais:

\textbf{(1) Kruskal não tem consciência da raiz.} O algoritmo de Kruskal seleciona arestas apenas pelo menor custo, sem nenhuma noção de qual vértice é a raiz nem de quais sub-árvores se formam ao redor dela. Para o Kruskal, a aresta $(3, 7)$ é tão boa quanto a aresta $(0, 5)$ --- o único critério é o custo. O CMSTP, porém, exige que a capacidade seja respeitada por \emph{sub-árvores enraizadas em $r$}, o que requer que o algoritmo saiba a posição da raiz a todo instante.

\textbf{(2) A AGM gerada pelo Kruskal quase sempre viola a capacidade.} A AGM ótima sem restrições frequentemente agrupa muitos terminais de alta demanda em um mesmo ramo, pois isso é mais barato em termos de arestas. Ao verificar a restrição $\sum_{i \in S} d_i \leq Q$ sobre essa árvore, muitos ramos a violam. Corrigir a violação exigiria cortar ramos e reconectá-los à raiz, aumentando o custo e resolvendo, na prática, um novo problema combinatório.

Precisamos, portanto, de um algoritmo que: (a) saiba quem é a raiz; (b) controle a demanda por componente durante a construção; e (c) meça a vantagem de uma fusão \emph{em relação ao custo de conexão com a raiz}. É exatamente isso que a heurística de Esau--Williams faz.


%=======================================================
%                 A L G O R I T M O S
%=======================================================
\section{Algoritmos Implementados}
\label{secAlgoritmos}

Nesta seção são detalhados os três algoritmos propostos. Todos compartilham o mesmo mecanismo de construção baseado na heurística de Esau--Williams; o que os diferencia é a estratégia de seleção do candidato a ser fundido.

\subsection{Representação da solução}

A solução é representada por um \textbf{vetor de pais}: \texttt{pai[v]} é o predecessor de $v$ no caminho até a raiz. Essa estrutura garante que o resultado seja sempre uma árvore (cada vértice tem exatamente um pai) e facilita o cálculo do custo total ($\sum_v c(v, \texttt{pai}[v])$) e a verificação de viabilidade (subindo de cada terminal até sua sub-raiz e acumulando a demanda por s-tree).

\subsection{Algoritmo Guloso}
\label{subSecAlgGuloso}

O algoritmo guloso implementa a heurística de Esau--Williams de forma determinística. A ideia central é partir de uma solução cara (todos os terminais ligados diretamente à raiz) e ir \textbf{economizando} iterativamente: a cada passo, funde-se a par de componentes que oferece a maior \textbf{economia}, entendida como a diferença entre o custo do \emph{portão} (gate) da componente --- a aresta mais barata conectando a componente à raiz --- e o custo da aresta interna que une as duas componentes. A fusão só é realizada se a demanda total resultante não ultrapassar $Q$.

A relação com o Kruskal fica clara na estrutura de dados: ambos utilizam \textit{union-find} para rastrear e fundir componentes. A diferença é que Kruskal seleciona pelo menor custo de aresta, enquanto Esau--Williams seleciona pela maior economia relativa ao portão, respeitando a restrição de capacidade.

O Algoritmo~\ref{algGuloso} descreve a construção. Na linha~1, cada terminal é inicializado como sua própria componente, com portão igual à aresta direta para a raiz. Nas linhas~4--10, para cada terminal~$i$ busca-se o terminal~$j$ mais próximo em outra componente cuja fusão não viole~$Q$; a economia de fundir $C_i$ com $C_j$ pela aresta $(i,j)$ é calculada na linha~9. Nas linhas~11--12, escolhe-se o par de maior economia positiva; se não houver, o algoritmo termina. As linhas~13--15 realizam a fusão e atualizam o portão da nova componente. Ao final (linhas~17--18), cada componente conecta-se à raiz pelo seu portão e uma BFS orienta a árvore.

\begin{algorithm}[H]
\caption{Esau--Williams Guloso}
\label{algGuloso}
\KwIn{Grafo $G=(V,E)$, raiz $r$, capacidade $Q$}
\KwOut{Árvore geradora $T$ enraizada em $r$ com demanda $\leq Q$ por s-tree}

\ForEach{$v \in V \setminus \{r\}$}{
    $\text{comp}[v] \leftarrow v$;\quad $\text{dem}[v] \leftarrow d_v$;\quad $\text{gate}[v] \leftarrow c(r,v)$\;
}
\Repeat{\textnormal{nenhum candidato com economia $> 0$}}{
    candidatos $\leftarrow \emptyset$\;
    \ForEach{$i \in V \setminus \{r\}$}{
        $c_i \leftarrow \text{find}(i)$\;
        $j^* \leftarrow \arg\min_{j:\,\text{find}(j)\neq c_i,\;\text{dem}[c_i]+\text{dem}[\text{find}(j)]\leq Q} c(i,j)$\;
        \If{$j^*$ existe}{
            $\text{economia}(i) \leftarrow \text{gate}[c_i] - c(i,j^*)$\;
            \If{$\text{economia}(i) > 0$}{
                adiciona $(i,j^*,\text{economia}(i))$ a candidatos\;
            }
        }
    }
    \If{\textnormal{candidatos} $= \emptyset$}{\textbf{break}\;}
    $(i^*, j^*) \leftarrow$ par com maior economia em candidatos\;
    adiciona aresta $(i^*,j^*)$ a $T$\;
    $c_i \leftarrow \text{find}(i^*)$;\quad $c_j \leftarrow \text{find}(j^*)$\;
    $\text{union}(c_i, c_j)$;\quad $\text{dem}[\cdot] \leftarrow \text{dem}[c_i]+\text{dem}[c_j]$\;
    $\text{gate}[\cdot] \leftarrow \min(\text{gate}[c_i], \text{gate}[c_j])$\;
}
\ForEach{\textnormal{componente } $C$}{
    adiciona aresta $(\text{gate\_node}[C],\, r)$ a $T$\;
}
Orienta $T$ por BFS a partir de $r$ (preenche vetor de pais)\;
\Return{$T$}
\end{algorithm}

Por ser determinístico, o algoritmo guloso sempre produz a mesma solução para uma dada instância e valor de $Q$, servindo como \textit{baseline} de comparação.

\subsection{Algoritmo Guloso Randomizado}
\label{subSecAlgRand}

A principal limitação do algoritmo guloso é sua natureza determinística: ao escolher sempre o melhor candidato, ele pode convergir para um mínimo local e ignorar outras regiões do espaço de soluções que poderiam produzir árvores de menor custo. Para superar essa limitação, adotou-se a metaheurística GRASP (\textit{Greedy Randomized Adaptive Search Procedure}) \cite{feo1995}.

Em vez de selecionar sempre o candidato de maior economia, o algoritmo constrói a \textbf{Lista Restrita de Candidatos (RCL)}, composta por todos os pares cuja economia esteja dentro de um limiar controlado pelo parâmetro $\alpha \in [0,1]$:

\begin{equation}
    \text{RCL} = \big\{\, i : \text{economia}(i) \geq e_{\max} - \alpha \cdot (e_{\max} - e_{\min}) \,\big\},
\end{equation}

e sorteia uniformemente um candidato da RCL. Quando $\alpha = 0$, apenas o melhor candidato integra a RCL e o algoritmo recai no guloso puro; quando $\alpha = 1$, todos os candidatos são elegíveis e a escolha é totalmente aleatória. Valores intermediários equilibram a intensificação (explorar a vizinhança da melhor solução) e a diversificação (explorar novas regiões). Foram testados os valores $\alpha \in \{0{,}1;\; 0{,}2;\; 0{,}3\}$.

A construção é repetida \texttt{iteracoes} vezes (mínimo 30, conforme a especificação) e a melhor solução encontrada é retornada, conforme o Algoritmo~\ref{algRand}.

\begin{algorithm}[H]
\caption{Esau--Williams Guloso Randomizado (GRASP)}
\label{algRand}
\KwIn{Grafo $G$, raiz $r$, capacidade $Q$, parâmetros $\alpha$ e \textit{iter}}
\KwOut{Melhor árvore $T^*$ encontrada nas \textit{iter} construções}

$T^* \leftarrow \emptyset$;\quad $c^* \leftarrow +\infty$\;
\For{$t \leftarrow 1$ \KwTo \textit{iter}}{
    $T \leftarrow$ \textsc{ConstruçãoRCL}$(G, r, Q, \alpha)$ \tcp*{usa RCL em vez do melhor}
    \If{$\text{custo}(T) < c^*$}{
        $T^* \leftarrow T$;\quad $c^* \leftarrow \text{custo}(T)$\;
    }
}
\Return{$T^*$}
\end{algorithm}

\subsection{Algoritmo Guloso Randomizado Reativo}
\label{subSecAlgRandReact}

A escolha do parâmetro $\alpha$ no algoritmo randomizado é crítica: um $\alpha$ muito pequeno limita a diversificação, enquanto um $\alpha$ muito grande prejudica a qualidade das soluções. O algoritmo randomizado reativo \cite{prais2000} elimina essa necessidade de ajuste manual ao \textbf{aprender} automaticamente quais valores de $\alpha$ funcionam melhor para cada instância.

O algoritmo mantém um conjunto de $m$ valores $\{\alpha_1, \ldots, \alpha_m\}$, cada um com uma probabilidade $p_k$ de ser sorteado. As probabilidades são atualizadas a cada \textbf{bloco} de iterações (tamanho 30--50) pela regra de Prais \& Ribeiro:

\begin{equation}
    q_k = \left(\frac{c^*}{\bar{c}_k}\right)^{\delta}, \qquad p_k = \frac{q_k}{\displaystyle\sum_{j=1}^{m} q_j},
\end{equation}

onde $c^*$ é o melhor custo encontrado até o momento, $\bar{c}_k$ é o custo médio das soluções geradas com $\alpha_k$ e $\delta = 10$ é o parâmetro de sensibilidade. Como $\bar{c}_k \geq c^*$, a razão $c^*/\bar{c}_k \in (0,1]$: valores de $\alpha$ que geram soluções de custo médio próximo ao melhor ganham peso maior, enquanto os de pior desempenho são penalizados --- mas não eliminados, para preservar a diversificação.

Foram utilizados os valores $\alpha \in \{0{,}05;\; 0{,}10;\; 0{,}15;\; 0{,}20;\; 0{,}30;\; 0{,}50\}$, com 300 iterações e blocos de 30 iterações, conforme a especificação do trabalho. O Algoritmo~\ref{algReativo} descreve o mecanismo reativo e a Figura~\ref{fig:fluxograma} apresenta o fluxo completo.

\begin{algorithm}[H]
\caption{Esau--Williams Guloso Randomizado Reativo}
\label{algReativo}
\KwIn{Grafo $G$, raiz $r$, capacidade $Q$, conjunto $\{\alpha_1,\ldots,\alpha_m\}$, parâmetros \textit{iter} e \textit{bloco}}
\KwOut{Melhor árvore $T^*$ e melhor $\alpha$ encontrado}

$p_k \leftarrow 1/m$ para todo $k$;\quad $T^* \leftarrow \emptyset$;\quad $c^* \leftarrow +\infty$\;
\For{$t \leftarrow 1$ \KwTo \textit{iter}}{
    Sorteia índice $k$ conforme a distribuição $\{p_k\}$\;
    $T \leftarrow$ \textsc{ConstruçãoRCL}$(G, r, Q, \alpha_k)$\;
    Registra custo $c(T)$ para $\alpha_k$\;
    \If{$c(T) < c^*$}{$T^* \leftarrow T$;\quad $c^* \leftarrow c(T)$;\quad $\alpha^* \leftarrow \alpha_k$\;}
    \If{$t \bmod \textit{bloco} = 0$}{
        \ForEach{$k$}{
            $q_k \leftarrow (c^*/\bar{c}_k)^{\delta}$ \tcp*{$\bar{c}_k$: custo médio de $\alpha_k$}
        }
        $p_k \leftarrow q_k / \sum_j q_j$ para todo $k$\;
    }
}
\Return{$T^*$, $\alpha^*$}
\end{algorithm}

\begin{figure}[ht]
\centering
\resizebox{0.55\textwidth}{!}{
    % Fluxograma do algoritmo guloso randomizado reativo
\usetikzlibrary{arrows.meta, shapes.geometric}
\tikzset{%
    >={Latex[width=2mm,length=2mm]},
    bloco/.style    = {rectangle, rounded corners, draw=black, fill=blue!10,
                       minimum width=4.5cm, minimum height=0.9cm,
                       text centered, font=\small},
    decisao/.style  = {diamond, draw=black, fill=orange!20, aspect=2.5,
                       text centered, font=\small},
    termino/.style  = {rectangle, rounded corners, draw=black, fill=gray!20,
                       minimum width=4.5cm, minimum height=0.9cm,
                       text centered, font=\small},
}
\begin{tikzpicture}[node distance=1.3cm, every node/.style={align=center}]
    \node (ini)   [termino]  {Início};
    \node (prob)  [bloco, below of=ini]   {Sorteia $\alpha_k$ conforme $\{p_k\}$};
    \node (cons)  [bloco, below of=prob]  {Constrói solução com Esau--Williams};
    \node (upd)   [bloco, below of=cons]  {Registra custo; atualiza melhor};
    \node (bloco) [decisao, below of=upd, yshift=-0.3cm] {Fim do bloco?};
    \node (prob2) [bloco, right of=bloco, xshift=3.5cm] {Reajusta $p_k$ (Prais \& Ribeiro)};
    \node (iter)  [decisao, below of=bloco, yshift=-0.3cm] {Fim das iterações?};
    \node (fim)   [termino, below of=iter, yshift=-0.3cm] {Retorna melhor solução};

    \draw[->] (ini)   -- (prob);
    \draw[->] (prob)  -- (cons);
    \draw[->] (cons)  -- (upd);
    \draw[->] (upd)   -- (bloco);
    \draw[->] (bloco) -- node[above, font=\small] {Sim} (prob2);
    \draw[->] (prob2) |- (iter);
    \draw[->] (bloco) -- node[left,  font=\small] {Não} (iter);
    \draw[->] (iter)  -- node[left,  font=\small] {Sim} (fim);
    \draw[->] (iter.west) -- ++(-2,0) |- node[left, font=\small, pos=0.25] {Não} (prob.west);
\end{tikzpicture}

}
\caption{Fluxograma do algoritmo guloso randomizado reativo.}
\label{fig:fluxograma}
\end{figure}


%=======================================================
%              E X P E R I M E N T O S
%=======================================================
\section{Experimentos computacionais}
\label{secResultados}

\subsection{Descrição das instâncias}
\label{subSecInstancias}

As instâncias utilizadas provêm da OR-Library \cite{beasley1990}, conjunto de referência amplamente utilizado na literatura do CMSTP. Utilizamos dois arquivos: \texttt{capmst1.txt} (grupos \texttt{tc80} e \texttt{te80}) e \texttt{capmst2.txt} (grupos \texttt{cm50}, \texttt{cm100} e \texttt{cm200}), lidos diretamente no formato nativo da OR-Library. Cada grupo contém 5 instâncias, totalizando 75 configurações distintas (5 grupos $\times$ 5 instâncias $\times$ 3 valores de $Q$).

\begin{table}[!h]
\centering
\caption{Grupos de instâncias utilizados nos experimentos}
\begin{tabular}{|l|c|c|c|c|}
\hline
\multicolumn{1}{|c|}{\textbf{Grupo}} & \textbf{Terminais} & \textbf{Demandas} & \textbf{Capacidades $Q$} & \textbf{Instâncias} \\ \hline
\texttt{tc80}  & 80 & Unitárias      & 5, 10, 20       & \texttt{tc80-1.txt} a \texttt{tc80-5.txt}   \\ \hline
\texttt{te80}  & 80 & Unitárias      & 5, 10, 20       & \texttt{te80-1.txt} a \texttt{te80-5.txt}   \\ \hline
\texttt{cm50}  & 50 & Não unitárias  & 200, 400, 800   & \texttt{cm50r1.dat} a \texttt{cm50r5.dat}   \\ \hline
\texttt{cm100} & 100 & Não unitárias & 200, 400, 800   & \texttt{cm100r1.dat} a \texttt{cm100r5.dat} \\ \hline
\texttt{cm200} & 200 & Não unitárias & 200, 400, 800   & \texttt{cm200r1.dat} a \texttt{cm200r5.dat} \\ \hline
\end{tabular}
\label{tabInstancias}
\end{table}

Os grupos \texttt{tc80} e \texttt{te80} possuem demandas unitárias ($d_i = 1$ para todo terminal), enquanto os grupos \texttt{cm} possuem demandas variáveis lidas do bloco \texttt{priz} do arquivo OR-Library. A capacidade $Q$ não está nos arquivos de instância; é passada como parâmetro de linha de comando, pois na literatura o mesmo grafo é testado com diferentes valores de $Q$.

\subsection{Ambiente computacional e parâmetros}
\label{subSecAmbiente}

O código foi implementado em C++ (padrão C++17) e compilado com \texttt{g++} com otimização \texttt{-O2}. Os experimentos foram executados em ambiente Linux. O gerador de números pseudoaleatórios utilizado é o Mersenne Twister (\texttt{std::mt19937}), com semente inicial baseada no relógio de alta resolução, impressa ao início de cada execução para permitir reprodutibilidade.

Cada algoritmo foi executado \textbf{10 vezes} por instância, com sementes de 1 a 10. Os parâmetros adotados foram:

\begin{itemize}[noitemsep]
    \item \textbf{Guloso:} determinístico, uma única execução por instância.
    \item \textbf{Randomizado:} $\alpha \in \{0{,}1;\; 0{,}2;\; 0{,}3\}$; mínimo de \textbf{30 construções} por execução.
    \item \textbf{Reativo:} $\alpha \in \{0{,}05;\; 0{,}10;\; 0{,}15;\; 0{,}20;\; 0{,}30;\; 0{,}50\}$; \textbf{300 construções} por execução; blocos de \textbf{30 iterações}.
\end{itemize}

\subsection{Resultados obtidos}
\label{subSecResultados}

O desvio percentual de cada resultado em relação à melhor solução conhecida da literatura \cite{ruiz2015} é calculado como:
\begin{equation}
    \text{desvio}(\%) = \frac{c_{\text{heurística}} - c^*}{c^*} \times 100.
\end{equation}

A Tabela~\ref{tabResultMelhor} apresenta o desvio percentual da \textbf{melhor} solução obtida (nas 10 execuções para os algoritmos randomizados) para cada grupo de instâncias. As células em negrito indicam o melhor resultado por linha.

\begin{table}[htbp]
\centering
\caption{Desvio percentual médio da \textbf{melhor} solução alcançada por cada algoritmo em relação à melhor solução conhecida, por grupo de instâncias}
\renewcommand{\arraystretch}{1.2}
{
\footnotesize
\begin{tabular}{|c|c|c|c|c|c|c|}
\hline
\textbf{Grupo} & \textbf{Melhor} & \textbf{Guloso} & \multicolumn{3}{c|}{\textbf{Randomizado}} & \textbf{Reativo} \\ \cline{4-6}
 & \textbf{conhecida} & & \multicolumn{1}{c|}{\textbf{0,10}} & \multicolumn{1}{c|}{\textbf{0,20}} & \multicolumn{1}{c|}{\textbf{0,30}} & \\
\hline
tc80  & -- &  &  &  &  &  \\
te80  & -- &  &  &  &  &  \\
cm50  & -- &  &  &  &  &  \\
cm100 & -- &  &  &  &  &  \\
cm200 & -- &  &  &  &  &  \\
\hline
\textbf{Média} &  &  &  &  &  &  \\
\hline
\end{tabular}
}
\label{tabResultMelhor}
\end{table}

A Tabela~\ref{tabResultMedia} apresenta o desvio percentual da \textbf{média} das 10 execuções de cada algoritmo, medida que reflete a consistência do método ao longo de múltiplas execuções com sementes diferentes.

\begin{table}[htbp]
\centering
\caption{Desvio percentual médio da \textbf{média das 10 execuções} de cada algoritmo em relação à melhor solução conhecida, por grupo de instâncias}
\renewcommand{\arraystretch}{1.2}
{
\footnotesize
\begin{tabular}{|c|c|c|c|c|c|c|}
\hline
\textbf{Grupo} & \textbf{Melhor} & \textbf{Guloso} & \multicolumn{3}{c|}{\textbf{Randomizado}} & \textbf{Reativo} \\ \cline{4-6}
 & \textbf{conhecida} & & \multicolumn{1}{c|}{\textbf{0,10}} & \multicolumn{1}{c|}{\textbf{0,20}} & \multicolumn{1}{c|}{\textbf{0,30}} & \\
\hline
tc80  & -- &  &  &  &  &  \\
te80  & -- &  &  &  &  &  \\
cm50  & -- &  &  &  &  &  \\
cm100 & -- &  &  &  &  &  \\
cm200 & -- &  &  &  &  &  \\
\hline
\textbf{Média} &  &  &  &  &  &  \\
\hline
\end{tabular}
}
\label{tabResultMedia}
\end{table}

A Tabela~\ref{tabResultTempo} apresenta o tempo médio de execução (em segundos) de cada algoritmo por grupo de instâncias.

\begin{table}[htbp]
\centering
\caption{Tempo médio de execução de cada algoritmo (em segundos), por grupo de instâncias}
\renewcommand{\arraystretch}{1.2}
{
\footnotesize
\begin{tabular}{|c|c|c|c|c|c|}
\hline
\textbf{Grupo} & \textbf{Guloso} & \multicolumn{3}{c|}{\textbf{Randomizado}} & \textbf{Reativo} \\ \cline{3-5}
 & & \multicolumn{1}{c|}{\textbf{0,10}} & \multicolumn{1}{c|}{\textbf{0,20}} & \multicolumn{1}{c|}{\textbf{0,30}} & \\
\hline
tc80  &  &  &  &  &  \\
te80  &  &  &  &  &  \\
cm50  &  &  &  &  &  \\
cm100 &  &  &  &  &  \\
cm200 &  &  &  &  &  \\
\hline
\textbf{Média} &  &  &  &  &  \\
\hline
\end{tabular}
}
\label{tabResultTempo}
\end{table}

\colorbox{pink}{
\begin{minipage}[c]{13cm}
[\textbf{Nota}:] As tabelas acima devem ser preenchidas com os resultados reais após a execução do script \texttt{scripts/experimentos.py}. A análise deve destacar quais algoritmos foram melhores, se o tamanho da instância ou o valor de $Q$ influencia os resultados, e comparar o comportamento do reativo com os alphas fixos.
\end{minipage}}


%=======================================================
%                 C O N C L U S Õ E S
%=======================================================
\section{Conclusões}
\label{secConclusoes}

Neste trabalho, implementamos três heurísticas construtivas para o CMSTP: um algoritmo guloso determinístico, um guloso randomizado (GRASP) e um guloso randomizado reativo. Todos os três compartilham o mecanismo de construção de Esau--Williams, heurística clássica do problema que parte de uma solução cara (todos os terminais conectados diretamente à raiz) e vai realizando fusões de componentes com base no conceito de \textbf{economia}, sempre respeitando a restrição de capacidade.

A escolha do Esau--Williams em detrimento de uma adaptação direta do Kruskal é justificada por razões fundamentais: o Kruskal não distingue a raiz dos demais vértices, não mantém controle de demanda por componente e pode gerar soluções completamente inviáveis. O Esau--Williams, por sua vez, foi projetado especificamente para o CMSTP, garantindo viabilidade por construção. Apesar de ambos utilizarem \textit{union-find}, seus critérios de seleção são essencialmente diferentes: menor custo de aresta (Kruskal) versus maior economia relativa ao portão da componente (Esau--Williams).

Entre as principais dificuldades encontradas destacam-se a leitura do formato nativo da OR-Library (que embute múltiplas instâncias em um único arquivo e remapeia o depósito para o último nó) e a calibração do parâmetro $\alpha$ do GRASP, motivação direta para a implementação do algoritmo reativo.

Como trabalhos futuros, a adição de uma fase de \textbf{busca local} após a construção --- por exemplo, mover um terminal de uma s-tree para outra quando isso reduzir o custo sem violar $Q$ (\textit{node insertion}) --- poderia reduzir significativamente o gap em relação às melhores soluções conhecidas.


%=======================================================
%             R E F E R Ê N C I A S
%=======================================================
\bibliographystyle{plain}
\bibliography{bibliografia}

\end{document}
